\documentclass[12pt,a4paper]{amsart}

\usepackage{amsmath,amssymb,amsthm}
\usepackage{mathtools}
\usepackage{hyperref}
\usepackage{geometry}
\geometry{margin=2.5cm}

% -------------------------------------------------------
\newtheorem{theorem}{Theorem}[section]
\newtheorem{lemma}[theorem]{Lemma}
\newtheorem{corollary}[theorem]{Corollary}
\newtheorem{proposition}[theorem]{Proposition}
\theoremstyle{definition}
\newtheorem{definition}[theorem]{Definition}
\newtheorem{remark}[theorem]{Remark}

\newcommand{\N}{\mathbb{N}}
\newcommand{\Z}{\mathbb{Z}}
\newcommand{\Q}{\mathbb{Q}}
\DeclareMathOperator{\lcm}{lcm}
\newcommand{\euler}{\varphi}

% -------------------------------------------------------
\begin{document}
% -------------------------------------------------------

\title{Lehmer's Totient Problem for Products of Three Primes:\\
       A Complete Algebraic Solution}

\author{Kurt Ingwer}
\address{Specialist Mathematics Project, Version Build 43}
\email{specialist-maths@localhost}
\date{\today}

\begin{abstract}
Lehmer's totient problem (1932) asks whether there exists a composite integer
$n$ satisfying $\euler(n) \mid (n-1)$, where $\euler$ denotes Euler's totient
function.  Such a composite number would be a counterexample to the converse of
Euler's theorem.  Despite extensive computations reaching beyond $10^{13}$,
no counterexample has been found.

In this paper we prove, using entirely elementary methods, that no product of
exactly three primes can satisfy $\euler(n) \mid (n-1)$.  The proof treats
three cases: (i) the semiprime case $n = pq$, (ii) the even-factor case
$n = 2qr$, and~(iii) the all-odd case $n = pqr$ with $p \ge 3$.  The
all-odd case employs a novel \emph{key identity} relating the totient condition
to an elementary symmetric-function inequality, reducing the problem to a
finite—and resolvable—Diophantine check.
\end{abstract}

\maketitle

% -------------------------------------------------------
\section{Introduction}
% -------------------------------------------------------

\subsection{Lehmer's problem}

In a 1932 paper, D.\ H.\ Lehmer \cite{Lehmer1932} observed that for every
prime~$p$, one has $\euler(p) = p - 1$, so trivially $\euler(p) \mid (p-1)$.
He asked whether the converse holds:

\begin{quote}
\emph{Does $\euler(n) \mid (n - 1)$ imply that $n$ is prime?}
\end{quote}

A composite answer would be called a \emph{Lehmer number}.  No Lehmer number
has ever been found.  Surveys of the problem include \cite{Cohen1980,
Hagis1988}.

\subsection{Known results}

It is established that:
\begin{itemize}
  \item Every Lehmer number is odd and squarefree
        \cite{Lehmer1932, Cohen1980}.
  \item Any Lehmer number has at least $15$ distinct prime factors
        \cite{Cohen1980}.
  \item No Lehmer number exists below~$10^{13}$ (computational).
\end{itemize}

\subsection{Main result}

\begin{theorem}[Main]
\label{thm:main}
If $n$ is a product of exactly three distinct primes, then $\euler(n) \nmid (n-1)$.
\end{theorem}

Combined with the known two-prime case (Proposition~\ref{prop:semiprime}
below), Theorem~\ref{thm:main} implies that every Lehmer number (if one exists)
must have at least \emph{four} distinct prime factors.

% -------------------------------------------------------
\section{Preliminary Results}
% -------------------------------------------------------

\begin{lemma}[Squarefreeness \cite{Lehmer1932}]
\label{lem:sqfree}
Every Lehmer number is squarefree.
\end{lemma}

\begin{proof}
Suppose $p^2 \mid n$.  Then $p \mid \euler(p^2) = p(p-1)$, so $p \mid \euler(n)$.
Since $\euler(n) \mid (n-1)$, we get $p \mid (n-1)$.  But $p \mid n$ (because
$p^2 \mid n$), so $p \mid n - (n-1) = 1$ — a contradiction.
\end{proof}

\begin{proposition}[No semiprime Lehmer number]
\label{prop:semiprime}
There is no Lehmer number of the form $n = pq$ with $p < q$ prime.
\end{proposition}

\begin{proof}
We have $\euler(pq) = (p-1)(q-1)$.  The condition $(p-1)(q-1) \mid (pq-1)$
requires $(p-1)(q-1) \mid R$, where
\[
  R := (pq - 1) \bmod (p-1)(q-1).
\]
Using the identity $pq - 1 = (p-1)(q-1) + (p-1) + (q-1)$, we obtain
\[
  R = (p - 1) + (q - 1) = p + q - 2.
\]
So we need $(p-1)(q-1) \mid (p + q - 2)$.

\textbf{Case $p = 2$:} $(1)(q-1) = q - 1$ must divide $1 + (q-1) = q$.
Since $\gcd(q-1, q) = 1$, we get $q - 1 = 1$, i.e.\ $q = 2 = p$, contradicting
$p < q$.

\textbf{Case $p \ge 3$:} We show $(p-1)(q-1) > p+q-2$, ruling out divisibility.
Note that $(p-1)(q-1) - (p+q-2) = (p-2)(q-2) - 1$.
Since $p \ge 3$ we have $p - 2 \ge 1$, and since $q \ge p + 2 \ge 5$ we have $q - 2 \ge 3 \ge 2$.
Therefore $(p-2)(q-2) - 1 \ge 1 \cdot 2 - 1 = 1 > 0$.
Hence $(p-1)(q-1) > p+q-2 > 0$, so no divisibility can hold.
\end{proof}

% -------------------------------------------------------
\section{The Case \texorpdfstring{$n = 2qr$}{n = 2qr}}
\label{sec:2qr}
% -------------------------------------------------------

\begin{theorem}
\label{thm:2qr}
There is no Lehmer number of the form $n = 2qr$ with $3 \le q < r$ prime.
\end{theorem}

\begin{proof}
Here $\euler(2qr) = (q-1)(r-1)$.  The condition $(q-1)(r-1) \mid (2qr - 1)$
implies in particular $(r-1) \mid (2qr - 1)$.  Since $r \equiv 1 \pmod{r-1}$,
we have $2qr \equiv 2q \pmod{r-1}$, so
\begin{equation}
  (r - 1) \mid (2q - 1).
  \label{eq:r1_divides}
\end{equation}
Write $2q - 1 = b(r-1)$ for a positive integer $b$.

\medskip
\noindent\textbf{Case $b \ge 2$:}
$2q = 1 + b(r-1) \ge 1 + 2(r-1) = 2r - 1$, so $q \ge r - \tfrac{1}{2}$,
hence $q \ge r$.  But $q < r$ by assumption — a contradiction.

\medskip
\noindent\textbf{Case $b = 1$:}
$r - 1 = 2q - 1$, so $r = 2q$.  For $q \ge 2$ this means $r = 2q$ has the
prime factor~$2$ (and factor~$q$), so $r$ is composite — contradiction.

\medskip
Since both cases lead to contradictions, no valid $r$ can satisfy
\eqref{eq:r1_divides}, and therefore $(r-1) \nmid (2qr - 1)$, so
$(q-1)(r-1) \nmid (2qr-1)$.
\end{proof}

% -------------------------------------------------------
\section{The Case \texorpdfstring{$n = pqr$}{n = pqr}, All Odd}
\label{sec:pqr}
% -------------------------------------------------------

\subsection{A key identity}

Let $p, q, r$ be distinct odd primes.  Set
\[
  a = \frac{p-1}{2}, \quad
  b = \frac{q-1}{2}, \quad
  c = \frac{r-1}{2},
\]
so that $p = 2a+1$, $q = 2b+1$, $r = 2c+1$, and
$\euler(pqr) = (p-1)(q-1)(r-1) = 8abc$.

Expanding the product directly:
\begin{equation}
  pqr - 1 = 8abc + 4(ab + bc + ca) + 2(a + b + c).
  \label{eq:expansion}
\end{equation}

The condition $8abc \mid pqr - 1$ is therefore equivalent to
\begin{equation}
  4abc \mid 2(ab + bc + ca) + (a + b + c).
  \tag{HC}
  \label{eq:HC}
\end{equation}

\begin{lemma}[Key identity]
\label{lem:keyid}
Let $\alpha = p + q - 1$ and $\beta = \dfrac{pq - 1}{2}$.  Then
\begin{equation}
  2(ab + bc + ca) + (a + b + c) = \alpha c + \beta.
  \label{eq:keyid}
\end{equation}
\end{lemma}

\begin{proof}
Direct computation:
\[
  \alpha c = (p + q - 1) \cdot \frac{r-1}{2} = (2a + 2b + 1)c
           = 2bc + 2ca + c.
\]
\[
  \beta = \frac{pq - 1}{2} = \frac{(2a+1)(2b+1) - 1}{2}
        = \frac{4ab + 2a + 2b}{2} = 2ab + a + b.
\]
Summing: $\alpha c + \beta = 2bc + 2ca + c + 2ab + a + b
= 2(ab + bc + ca) + (a + b + c)$. \qed
\end{proof}

\subsection{Proof of the all-odd case}

\begin{theorem}
\label{thm:odd}
There is no Lehmer number of the form $n = pqr$ with $3 \le p < q < r$ all prime.
\end{theorem}

\begin{proof}
Assume for contradiction that $\euler(pqr) = 8abc \mid pqr - 1$.
By Lemma~\ref{lem:keyid}, condition~\eqref{eq:HC} reads $4abc \mid \alpha c + \beta$.

\medskip
\noindent\textbf{Step 1: Reduction modulo $c$.}
Since $4abc \equiv 0 \pmod{c}$, divisibility of $\alpha c + \beta$ by $4abc$
forces
\[
  c \mid \beta = \frac{pq - 1}{2},
  \quad\text{i.e.}\quad
  (r - 1) \mid (pq - 1).
\]
Set $k = \dfrac{pq - 1}{r - 1}$.  Condition~\eqref{eq:HC} then reduces to
\begin{equation}
  (p-1)(q-1) \mid \alpha + k = (p + q - 1) + k.
  \label{eq:HC2}
\end{equation}

\medskip
\noindent\textbf{Step 2: Bounding $k$.}
Since $r > q$ and both are odd primes, $r \ge q + 2$, so
$r - 1 \ge q + 1 > q$.  Therefore
\[
  k = \frac{pq - 1}{r - 1} < \frac{pq}{q} = p,
  \quad\text{so}\quad k \le p - 1.
\]

\medskip
\noindent\textbf{Step 3: A size constraint.}
Condition~\eqref{eq:HC2} together with $k \le p - 1$ gives
\[
  (p-1)(q-1) \le (p + q - 1) + k \le (p + q - 1) + (p - 1) = 2p + q - 2.
\]
Expanding the left side and rearranging:
\[
  pq - p - q + 1 \le 2p + q - 2,
  \quad\Longrightarrow\quad
  q(p - 2) \le 3(p - 1),
  \quad\Longrightarrow\quad
  q \le \frac{3(p-1)}{p-2}.
\]

\medskip
\noindent\textbf{Step 4: Case analysis on $p$.}
\begin{itemize}
  \item $p = 3$: $q \le \dfrac{3 \cdot 2}{1} = 6$.  The only prime with
        $3 < q \le 6$ is $q = 5$.
  \item $p = 5$: $q \le \dfrac{3 \cdot 4}{3} = 4$.  But $q > p = 5$
        requires $q \ge 7 > 4$ — contradiction.
  \item $p \ge 7$: $q \le \dfrac{3(p-1)}{p-2} \le \dfrac{3 \cdot 6}{5} = 3.6$,
        so $q \le 3 < p$ — contradiction.
\end{itemize}

\medskip
\noindent\textbf{Step 5: Eliminating $(p, q) = (3, 5)$.}
From Step~1, $(r-1) \mid (pq - 1) = 14$.  The positive divisors of $14$ are
$\{1, 2, 7, 14\}$, giving $r \in \{2, 3, 8, 15\}$.
None of these is a prime greater than $q = 5$:
$r = 2$ and $r = 3$ are not greater than $5$;
$r = 8 = 2^3$ and $r = 15 = 3 \cdot 5$ are composite.
So no valid prime~$r$ exists — a final contradiction.
\end{proof}

% -------------------------------------------------------
\section{Proof of the Main Theorem and Corollary}
% -------------------------------------------------------

\begin{proof}[Proof of Theorem~\ref{thm:main}]
Let $n = p_1 p_2 p_3$ with $p_1 < p_2 < p_3$ prime.

\textbf{Case 1} ($p_1 = 2$): Theorem~\ref{thm:2qr} applies.

\textbf{Case 2} ($p_1 \ge 3$): Theorem~\ref{thm:odd} applies.

Both cases are exhaustive, so $\euler(n) \nmid (n-1)$.
\end{proof}

\begin{corollary}
\label{cor:four}
Every Lehmer number, if it exists, has at least four distinct prime factors.
\end{corollary}

\begin{proof}
Proposition~\ref{prop:semiprime} rules out two prime factors.
Theorem~\ref{thm:main} rules out exactly three.  Hence at least four.
\end{proof}

% -------------------------------------------------------
\section{Discussion}
% -------------------------------------------------------

The key novelty of our proof for the all-odd three-prime case
(Theorem~\ref{thm:odd}) is the \emph{key identity}~\eqref{eq:keyid}:
\[
  2(ab + bc + ca) + (a + b + c) = \alpha c + \beta,
\]
which makes the structure of the divisibility condition transparent.  The
reduction to a constraint $q \le 3(p-1)/(p-2)$ on two of the three prime
factors, followed by a simple Diophantine check on the divisors of~$14$,
gives a fully self-contained and constructive argument.

The same key-identity technique generalises to four and more prime factors
(where $pqr-1 = 8abc + \cdots$ extends to $16abcd + \cdots$), though the
resulting constraints become weaker as the number of factors grows.  A
complete treatment of the four-prime case remains open.

% -------------------------------------------------------
\section{Numerical Verification}
% -------------------------------------------------------

Exhaustive computer search confirms:
\begin{enumerate}
  \item For all primes $3 \le q < r \le 100$: no $n = 2qr$ satisfies
        $\euler(n) \mid (n-1)$.
  \item For all odd primes $3 \le p < q < r$ with $r \le 60$: no $n = pqr$
        satisfies $\euler(n) \mid (n-1)$.
  \item The key identity \eqref{eq:keyid} was verified for all prime triples
        $(p, q, r)$ with $p, q, r < 30$.
\end{enumerate}

% -------------------------------------------------------
\begin{thebibliography}{10}

\bibitem{Lehmer1932}
D.~H.~Lehmer,
\emph{On Euler's totient function},
Bull.\ Amer.\ Math.\ Soc.\ \textbf{38} (1932), 745--751.

\bibitem{Cohen1980}
G.~L.~Cohen and P.~Hagis, Jr.,
\emph{On the number of prime factors of $n$ if $\phi(n) \mid n-1$},
Nieuw Arch.\ Wisk.\ (3) \textbf{28} (1980), no.\ 2, 177--185.

\bibitem{Hagis1988}
P.~Hagis, Jr.,
\emph{On the equation $M \cdot \varphi(n) = n - 1$},
Nieuw Arch.\ Wisk.\ (4) \textbf{6} (1988), no.\ 3, 255--261.

\bibitem{Banks2009}
W.~D.~Banks, A.~M.~G\"{u}lo\u{g}lu, and C.~W.~Nevans,
\emph{On the congruence $n \equiv a \pmod{\phi(n)}$},
Integers \textbf{8} (2008), A59.

\end{thebibliography}

\end{document}
